\documentclass[a4paper]{article}

\usepackage[fontset=fandol]{ctex}
\usepackage{amsmath, amssymb}
\usepackage{graphicx}
\usepackage[margin=2.3cm]{geometry}
\usepackage[most]{tcolorbox}
\usepackage{etoolbox}
\usepackage{listings}
\usepackage{hyperref}
\usepackage{xurl}
\usepackage{booktabs}
\usepackage{longtable}
\usepackage{tabularx}
\usepackage{array}
\usepackage{subcaption}
\usepackage{float}
\usepackage{enumitem}
\usepackage{tikz}
\usetikzlibrary{arrows.meta, positioning, shapes.geometric}

\hypersetup{
    colorlinks=true,
    linkcolor=blue!50!black,
    urlcolor=blue!60!black,
    citecolor=blue!50!black
}

\newtcolorbox{findingbox}[1]{
    enhanced,
    colback=green!5!white,
    colframe=green!45!black,
    colbacktitle=green!45!black,
    coltitle=white,
    fonttitle=\bfseries,
    title=#1,
    sharp corners
}

\newtcolorbox{evidencebox}[1]{
    enhanced,
    colback=blue!4!white,
    colframe=blue!65!black,
    colbacktitle=blue!65!black,
    coltitle=white,
    fonttitle=\bfseries,
    title=#1,
    sharp corners
}

\newtcolorbox{riskbox}[1]{
    enhanced,
    colback=red!4!white,
    colframe=red!70!black,
    colbacktitle=red!70!black,
    coltitle=white,
    fonttitle=\bfseries,
    title=#1,
    sharp corners
}

\newtcolorbox{methodbox}[1]{
    enhanced,
    colback=yellow!9!white,
    colframe=yellow!60!black,
    colbacktitle=yellow!60!black,
    coltitle=black,
    fonttitle=\bfseries,
    title=#1,
    sharp corners
}

\lstset{
    language=Python,
    basicstyle=\ttfamily\small,
    keywordstyle=\color{blue},
    stringstyle=\color{red!60!black},
    commentstyle=\color{green!50!black},
    breaklines=true,
    frame=single,
    numbers=left,
    numberstyle=\tiny\color{gray},
    captionpos=b,
    extendedchars=false
}

\newcolumntype{Y}{>{\raggedright\arraybackslash}X}
\setlist[itemize]{leftmargin=1.8em}

\newcommand{\reporttitle}{bilibili-api-python 在 video-note-pipeline 中的集成研究}
\newcommand{\reportsubtitle}{面向 Bilibili 官方字幕 fallback 的实现、依赖、交互与持久化设计}
\newcommand{\reportauthors}{Codex}
\newcommand{\reportdate}{2026-03-26}
\newcommand{\researchquestion}{在保留现有 \texttt{yt-dlp} 主链路的前提下，\texttt{bilibili-api-python} 是否适合接入 \texttt{video-note-pipeline} 作为 Bilibili 专用 fallback，以及它应该如何安装、认证、调用、落盘与暴露给用户。}
\newcommand{\reportscope}{聚焦当前本机仓库 \path{/home/fanghaotian/Desktop/src/video-note-pipeline} 的 Bilibili 字幕探测与媒体探测流程；不覆盖 YouTube，不设计最终 change，只给出可执行的工程结论。}
\newcommand{\reportaudience}{准备为 \texttt{video-note-pipeline} 编写下一轮 change / propose / implementation 的开发者。}
\newcommand{\sourcewindow}{官方资料与本地实验访问日期均截至 2026-03-26（Asia/Shanghai）。}
\newcommand{\reportcoverpath}{}

\begin{document}

\begin{titlepage}
\centering
{\Large 深度研究报告\par}
\vspace{1.0cm}
{\huge\bfseries \reporttitle\par}
\vspace{0.5cm}
{\Large \reportsubtitle\par}
\vspace{0.9cm}
{\large \reportauthors\par}
\vspace{0.2cm}
{\large \reportdate\par}
\vspace{1.0cm}

\ifdefempty{\reportcoverpath}{
{\small 本报告未使用封面图，重点保留可审计的源码、文档与本地实验 artifact。\par}
}{
\includegraphics[width=0.84\textwidth,height=0.42\textheight,keepaspectratio]{\reportcoverpath}\par
}

\vfill
\begin{tcolorbox}[width=0.94\textwidth, colback=black!2!white, colframe=black!60, sharp corners]
\textbf{核心问题}：\researchquestion\par
\textbf{研究范围}：\reportscope\par
\textbf{目标读者}：\reportaudience\par
\textbf{证据窗口}：\sourcewindow\par
\end{tcolorbox}
\end{titlepage}

\tableofcontents
\newpage

\section{执行摘要}

\begin{findingbox}{核心结论}
\texttt{bilibili-api-python} 适合接入当前项目，但最佳位置不是替换 \texttt{yt-dlp}，而是作为 \textbf{Bilibili 专用、带 cookies 的第二阶段 fallback}。理由有三点：
\begin{itemize}
\item 它能稳定拿到当前 \texttt{yt-dlp} 路径之外的 Bilibili 原生字幕元信息，关键接口是 \texttt{Video.get\_subtitle()} / \texttt{Video.get\_player\_info()}，并且本地已实测拿到 \texttt{ai-zh} 字幕。
\item 它可以直接复用我们现有的浏览器导出 cookies 文件，不需要引入新的登录交互；通过 \texttt{Credential.from\_cookies(...)} 即可桥接。
\item 它是全异步、依赖更重、且仓库许可证为 GPLv3+，因此作为全局主链路替代方案的代价偏高；放在 Bilibili fallback 层，收益和复杂度最平衡。
\end{itemize}
\end{findingbox}

\begin{findingbox}{推荐工程策略}
面向下一轮代码实现，建议采用以下策略：
\begin{itemize}
\item 保留当前 \texttt{yt-dlp} Bilibili 字幕路径作为第一层。
\item 当 \texttt{yt-dlp} 在 cookies 路径下仍未选中官方字幕时，且本地已有有效 Bilibili cookies，则进入 \texttt{bilibili-api-python} fallback。
\item fallback 内部复用已有 cookie 文件，调用 \texttt{get\_pages()}、\texttt{get\_cid()}、\texttt{get\_subtitle()}，然后对返回的 \texttt{subtitle\_url} 再发一次 HTTP GET，落盘原始 JSON，并归一化成当前项目已有的 transcript payload。
\item 用户侧不新增 QR 登录、短信登录、密码登录命令，继续统一使用 \texttt{video-note cookies-export bilibili}。
\end{itemize}
\end{findingbox}

\begin{riskbox}{本轮调研确认的主要风险}
\begin{itemize}
\item \textbf{许可证风险}：官方 README 明示 GPLv3+。如果 \texttt{video-note-pipeline} 会以依赖形式分发，必须先做许可证评估。
\item \textbf{刷新风险}：\texttt{credential.refresh()} 需要 \texttt{ac\_time\_value}。我们当前导出的 Netscape cookies 文件没有这个字段，因此不能直接依赖库内刷新逻辑。
\item \textbf{异步边界}：该库“全部是异步操作”。当前 CLI 以同步风格组织，因此需要一个边界层来管理事件循环与超时。
\item \textbf{原始字幕不是最终 transcript}：\texttt{get\_subtitle()} 返回的是字幕元信息，不是项目可直接消费的标准 transcript，仍需二次请求与归一化。
\end{itemize}
\end{riskbox}

\section{研究范围与方法}

\begin{methodbox}{研究方法}
本报告没有只读文档就下结论，而是把官方资料与本地可运行实验绑定在一起：
\begin{itemize}
\item 在线核验官方 README、PyPI、官方文档入口，确认版本、安装方式、请求客户端和登录模型。
\item 在隔离环境内真实安装 \texttt{bilibili-api-python==17.4.1}，并额外安装 \texttt{httpx} 与 \texttt{curl\_cffi} 做可运行验证。
\item 直接复用当前机器上的 \path{/home/fanghaotian/.config/video-note-pipeline/cookies/bilibili.txt}，验证能否转成 \texttt{Credential} 并拿到目标视频的字幕。
\item 对照当前仓库里的 Bilibili adapter 与 subtitle probe 实现，分析应该如何接入、如何落盘、如何保持用户交互稳定。
\end{itemize}
\end{methodbox}

\subsection{证据清单}

\begin{longtable}{p{1.5cm}p{1.5cm}p{4.0cm}p{6.5cm}}
\toprule
ID & 模态 & 标题 / 标签 & 链接或路径 \\
\midrule
\endfirsthead
\toprule
ID & 模态 & 标题 / 标签 & 链接或路径 \\
\midrule
\endhead
S1 & Web & 官方 README & \url{https://github.com/Nemo2011/bilibili-api} \\
S2 & Web & 官方文档入口 & \url{https://nemo2011.github.io/bilibili-api/} \\
S3 & Web & PyPI 项目页 & \url{https://pypi.org/project/bilibili-api-python/} \\
S4 & Web/Repo & 获取 Credential 文档 & \url{https://github.com/Nemo2011/bilibili-api/blob/main/docs/get-credential.md} \\
S5 & Web/Repo & 刷新 cookies 文档 & \url{https://github.com/Nemo2011/bilibili-api/blob/main/docs/refresh_cookies.md} \\
S6 & Web/Repo & 请求客户端文档 & \url{https://github.com/Nemo2011/bilibili-api/blob/main/docs/request_client.md} \\
S7 & Web/Repo & 配置文档 & \url{https://github.com/Nemo2011/bilibili-api/blob/main/docs/configuration.md} \\
S8 & Web/Repo & Video 模块文档 & \url{https://github.com/Nemo2011/bilibili-api/blob/main/docs/modules/video.md} \\
S9 & Web/Repo & Login v2 模块文档 & \url{https://github.com/Nemo2011/bilibili-api/blob/main/docs/modules/login_v2.md} \\
S10 & Local code & 当前项目 Bilibili adapter & \path{/home/fanghaotian/Desktop/src/video-note-pipeline/src/video_note_pipeline/adapters/bilibili.py} \\
S11 & Local code & 当前 subtitle probe 组装逻辑 & \path{/home/fanghaotian/Desktop/src/video-note-pipeline/src/video_note_pipeline/probes/subtitles.py} \\
S12 & Local code & 当前项目 pyproject 依赖声明 & \path{/home/fanghaotian/Desktop/src/video-note-pipeline/pyproject.toml} \\
S13 & Local experiment & cookies 转 Credential 验证 & \path{/home/fanghaotian/Desktop/src/agent-basic-skill/reports/2026-03-26-bilibili-api-python-integration/artifacts/credential-check.json} \\
S14 & Local experiment & refresh 失败原因验证 & \path{/home/fanghaotian/Desktop/src/agent-basic-skill/reports/2026-03-26-bilibili-api-python-integration/artifacts/credential-refresh-attempt.json} \\
S15 & Local experiment & public vs auth 行为差异 & \path{/home/fanghaotian/Desktop/src/agent-basic-skill/reports/2026-03-26-bilibili-api-python-integration/artifacts/public-vs-auth-behavior.json} \\
S16 & Local experiment & 字幕元信息与原始 JSON 摘要 & \path{/home/fanghaotian/Desktop/src/agent-basic-skill/reports/2026-03-26-bilibili-api-python-integration/artifacts/subtitle-metadata-and-json-summary.json} \\
S17 & Local experiment & 原始字幕 JSON 样本 & \path{/home/fanghaotian/Desktop/src/agent-basic-skill/reports/2026-03-26-bilibili-api-python-integration/artifacts/subtitle-json-sample.json} \\
S18 & Local experiment & 下载地址能力摘要 & \path{/home/fanghaotian/Desktop/src/agent-basic-skill/reports/2026-03-26-bilibili-api-python-integration/artifacts/download-url-summary.json} \\
S19 & Local experiment & QR 登录能力摘要 & \path{/home/fanghaotian/Desktop/src/agent-basic-skill/reports/2026-03-26-bilibili-api-python-integration/artifacts/qrcode-login-summary.json} \\
S20 & Local experiment & \texttt{curl\_cffi} 实测字幕检查 & \path{/home/fanghaotian/Desktop/src/agent-basic-skill/reports/2026-03-26-bilibili-api-python-integration/artifacts/curl-cffi-subtitle-check.json} \\
\bottomrule
\end{longtable}

\section{当前项目上下文}

当前项目的 Bilibili 逻辑仍然是 \texttt{yt-dlp} 主导。以 2026-03-26 本地代码为准，Bilibili adapter 的几个关键特征如下：
\begin{itemize}
\item 字幕请求走 \texttt{yt-dlp}，并且语言请求现在已经扩大为 \texttt{all,-live\_chat,-danmaku}。
\item 语言归一化已经能识别 \texttt{ai-zh}、\texttt{auto-zh}、\texttt{zh-Hans}、\texttt{zh-Hant} 等标签，并将其落入统一选择逻辑。
\item subtitle probe 的输出已经保留 \texttt{raw\_tracks}、\texttt{normalized\_tracks}、\texttt{excluded\_tracks} 与 \texttt{selection\_trace}，具备较好的可观测性。
\item cookies 获取已通过 \texttt{video-note cookies-export bilibili} 固化，并以 Netscape cookie file 形式存放在用户配置目录。
\end{itemize}

\begin{evidencebox}{为什么还要调研 bilibili-api-python}
虽然前一轮已经修复了 Bilibili \texttt{yt-dlp} 对 \texttt{ai-zh} 的识别问题，但 \texttt{yt-dlp} 本质上仍是站外抽象层。它能拿到 sidecar 字幕时很好用，但它并不暴露 Bilibili 原生的 \texttt{player\_info}、\texttt{need\_login\_subtitle}、\texttt{subtitle\_url} 等站内字段。对于“cookies 有效但 \texttt{yt-dlp} 仍未返回可用字幕”的 case，Bilibili 专用 fallback 仍然有价值。
\end{evidencebox}

\section{官方能力与安装模型}

\subsection{安装与依赖结构}

官方 README 明确给出安装方式：
\begin{itemize}
\item \texttt{pip install bilibili-api-python}
\item 另行安装一个异步请求库：\texttt{aiohttp}、\texttt{httpx} 或 \texttt{curl\_cffi}
\end{itemize}

本地隔离环境的实测安装版本是 \texttt{bilibili-api-python==17.4.1}。直接元数据依赖如下：
\begin{itemize}
\item \texttt{beautifulsoup4}
\item \texttt{colorama}
\item \texttt{lxml}
\item \texttt{pyyaml}
\item \texttt{brotli}
\item \texttt{qrcode}
\item \texttt{APScheduler}
\item \texttt{pillow}
\item \texttt{yarl}
\item \texttt{pycryptodomex}
\item \texttt{qrcode\_terminal}
\item 以及单独选择的 HTTP 客户端，例如 \texttt{httpx} 或 \texttt{curl\_cffi}
\end{itemize}

\begin{table}[H]
\centering
\begin{tabularx}{\textwidth}{p{3.0cm}p{3.1cm}Y}
\toprule
方案 & 需要安装 & 适合当前项目的原因 \\
\midrule
最小接入 & \texttt{bilibili-api-python} + \texttt{httpx} & 本地已验证能拿到目标视频的 \texttt{ai-zh} 字幕，依赖较轻，足够作为 fallback 起步。 \\
增强抗风控 & \texttt{bilibili-api-python} + \texttt{curl\_cffi} & 官方优先级更高，支持 impersonation。若后续遇到 412 或风控波动，可启用。 \\
不推荐现在就做 & 直接替换现有 \texttt{yt-dlp} Bilibili 主链路 & 需要引入更多异步边界、许可证评估和额外归一化代码，当前收益不如 fallback 明确。 \\
\bottomrule
\end{tabularx}
\caption{面向当前项目的安装策略}
\end{table}

\subsection{请求客户端与网络设置}

官方文档明确指出默认支持三类异步请求库，并按 \texttt{curl\_cffi > aiohttp > httpx} 的顺序选择。文档还给出如下可配置项：
\begin{itemize}
\item 通用：\texttt{proxy}、\texttt{timeout}、\texttt{verify\_ssl}、\texttt{trust\_env}
\item 库特有：\texttt{impersonate}、\texttt{http2}
\item Bilibili 特有：\texttt{wbi\_retry\_times}、\texttt{enable\_auto\_buvid}、\texttt{enable\_bili\_ticket}
\end{itemize}

\begin{findingbox}{对当前项目的建议}
\textbf{第一阶段默认用 \texttt{httpx}}。原因是：
\begin{itemize}
\item 我们已经在隔离环境里用 \texttt{httpx} 成功完成了 Bilibili 字幕元信息与字幕 JSON 的获取。
\item 它比 \texttt{curl\_cffi} 更容易维护，且 \texttt{video-note-pipeline} 当前不依赖 WebSocket。
\item 若后续遇到风控问题，再引入可选增强模式：\texttt{select\_client("curl\_cffi")} + \texttt{request\_settings.set("impersonate", "...") }。
\end{itemize}
\end{findingbox}

\subsection{许可证与分发问题}

官方 README 明确声明项目遵循 \textbf{GPLv3+}。这不是实现细节，而是设计边界：
\begin{itemize}
\item 如果 \texttt{video-note-pipeline} 只是本机工具或内部分发，依赖引入的工程阻力较小，但仍应在变更文档中显式记录。
\item 如果它将来要以闭源或混合许可方式对外分发，必须在引入前完成许可证评估。
\end{itemize}

\section{认证模型与用户交互}

\subsection{Credential 需要什么}

官方文档和源码都表明，\texttt{Credential} 可以从以下字段构造：
\begin{itemize}
\item \texttt{SESSDATA}
\item \texttt{bili\_jct}
\item \texttt{buvid3}
\item \texttt{buvid4}
\item \texttt{DedeUserID}
\item \texttt{ac\_time\_value}
\end{itemize}

其中 \texttt{ac\_time\_value} 只在 cookies 刷新时需要，来自浏览器 \texttt{localStorage}，并不在普通 cookie file 里。

\subsection{当前项目 cookies 能否直接复用}

答案是能，而且已经本地实测成功。当前 \texttt{video-note-pipeline} 导出的 Bilibili cookies 文件中，以下关键信息都能取到：
\begin{itemize}
\item \texttt{SESSDATA}
\item \texttt{bili\_jct}
\item \texttt{buvid3}
\item \texttt{buvid4}
\item \texttt{DedeUserID}
\end{itemize}

对应的本地 artifact 显示：
\begin{itemize}
\item \texttt{check\_valid = true}
\item \texttt{check\_refresh = true}
\item 但 \texttt{ac\_time\_value} 缺失
\end{itemize}

\begin{riskbox}{关于 refresh 的实际结论}
本地已经直接验证：当前导出的 cookie file \textbf{不能} 直接调用 \texttt{credential.refresh()}。错误是 \texttt{CredentialNoAcTimeValueException}。因此，若接入本库，推荐的策略不是“自动 refresh cookies”，而是继续沿用现有用户动作：
\begin{itemize}
\item cookie 失效时重新执行 \texttt{video-note cookies-export bilibili}
\item 不新增“自动刷新 Bilibili cookies”命令
\end{itemize}
\end{riskbox}

\subsection{是否要暴露 QR 登录 / 密码登录 / 短信登录}

官方库支持：
\begin{itemize}
\item \texttt{QrCodeLogin}
\item 密码登录
\item 短信登录
\end{itemize}

本地也已经验证 \texttt{QrCodeLogin} 的二维码生成路径存在。但对当前项目而言，这不是推荐方向：
\begin{itemize}
\item 我们已经有更稳定的浏览器 cookie 导出链路。
\item 新增登录交互会增加终端状态机、超时轮询、二维码展示与失败恢复复杂度。
\item 从用户心智上，统一“先在浏览器登录，再导出 cookies”更简单。
\end{itemize}

\begin{findingbox}{用户交互建议}
如果未来接入 \texttt{bilibili-api-python}，用户不应感知到新命令。推荐交互如下：
\begin{itemize}
\item 正常用户：继续只用 \texttt{video-note cookies-export bilibili --browser edge --json}
\item pipeline：自动检测已有 cookie file，并在 Bilibili fallback 中复用
\item 失败时：继续提示重导 cookies，而不是要求用户扫描二维码或粘贴 token
\end{itemize}
\end{findingbox}

\section{它能拿到什么信息}

\subsection{对字幕链路最关键的接口}

本次调研里最有价值的是以下几个视频接口：
\begin{itemize}
\item \texttt{Video.get\_info()}：公开视频基础信息
\item \texttt{Video.get\_pages()}：获取分 P 信息
\item \texttt{Video.get\_cid()}：定位具体页的 \texttt{cid}
\item \texttt{Video.get\_player\_info(cid=...)}：返回字幕、地区与最近播放信息，含字幕 JSON 链接
\item \texttt{Video.get\_subtitle(cid=...)}：从 \texttt{player\_info} 中抽出字幕区块
\item \texttt{Video.get\_download\_url(...)}：返回播放/下载地址信息
\end{itemize}

\subsection{公开访问与认证访问的实际边界}

对目标视频 \texttt{BV1DzALzQEHU} 的本地实验结果非常清楚：
\begin{itemize}
\item \texttt{get\_info()}、\texttt{get\_pages()}、\texttt{get\_cid()} 在匿名状态下可用。
\item 匿名调用 \texttt{get\_player\_info()} 会报 \texttt{CredentialNoSessdataException}。
\item 带当前 cookie file 转换得到的 \texttt{Credential} 后，\texttt{get\_subtitle()} 返回 1 条字幕轨，语言为 \texttt{ai-zh}。
\end{itemize}

\begin{table}[H]
\centering
\begin{tabularx}{\textwidth}{p{4.0cm}p{4.0cm}Y}
\toprule
能力 & 匿名访问 & 带当前 cookie file 的访问 \\
\midrule
\texttt{get\_info()} / \texttt{get\_pages()} / \texttt{get\_cid()} & 可用 & 可用 \\
\texttt{get\_player\_info()} / \texttt{get\_subtitle()} & 不可用，缺 \texttt{SESSDATA} & 可用，返回 \texttt{ai-zh} 字幕轨 \\
\texttt{get\_download\_url()} & 可用，但 DASH 视频流更少 & 可用，且本地实测看到更多 DASH 视频流 \\
\bottomrule
\end{tabularx}
\caption{Bilibili 目标视频的公开/认证能力对比}
\end{table}

\subsection{字幕接口真正返回什么}

\texttt{get\_subtitle()} 并不直接返回标准化后的 transcript，而是返回一个“字幕索引对象”。实测字段包括：
\begin{itemize}
\item \texttt{subtitles[].lan}，例如 \texttt{ai-zh}
\item \texttt{subtitles[].lan\_doc}，例如“中文”
\item \texttt{subtitles[].subtitle\_url}
\item \texttt{subtitles[].subtitle\_url\_v2}
\item \texttt{subtitles[].ai\_type} / \texttt{ai\_status}
\end{itemize}

其中 \texttt{subtitle\_url} 是协议相对地址，例如 \texttt{//aisubtitle...}。代码接入时必须先规范化为 \texttt{https://...} 再发起 HTTP GET。

拿到实际 JSON 后，顶层包含：
\begin{itemize}
\item \texttt{font\_size}
\item \texttt{font\_color}
\item \texttt{background\_alpha}
\item \texttt{background\_color}
\item \texttt{Stroke}
\item \texttt{type}
\item \texttt{lang}
\item \texttt{version}
\item \texttt{body[]}
\end{itemize}

其中真正要转成 transcript 的是 \texttt{body[]}：
\begin{itemize}
\item \texttt{from}：起始秒数
\item \texttt{to}：结束秒数
\item \texttt{content}：字幕文本
\item 其他补充字段如 \texttt{sid}、\texttt{location}、\texttt{music}
\end{itemize}

\begin{lstlisting}[caption={将现有 cookie file 转为 Credential 并获取字幕元信息的最小原型}]
import asyncio
from pathlib import Path
from bilibili_api import Credential, video, select_client


def load_netscape_cookies(cookie_file: Path) -> dict[str, str]:
    cookies = {}
    for raw in cookie_file.read_text(encoding="utf-8").splitlines():
        if not raw or raw.startswith("#"):
            continue
        parts = raw.split("\t")
        if len(parts) >= 7:
            cookies[parts[5]] = parts[6]
    return cookies


async def fetch_subtitle_index(bvid: str, cookie_file: Path) -> dict:
    select_client("httpx")
    credential = Credential.from_cookies(load_netscape_cookies(cookie_file))
    v = video.Video(bvid=bvid, credential=credential)
    cid = await v.get_cid(0)
    return await v.get_subtitle(cid=cid)


subtitle_index = asyncio.run(
    fetch_subtitle_index(
        "BV1DzALzQEHU",
        Path("/home/fanghaotian/.config/video-note-pipeline/cookies/bilibili.txt"),
    )
)
\end{lstlisting}

\section{本地实验结果}

\subsection{实验环境}

\begin{itemize}
\item 操作系统：Arch Linux（当前机器）
\item 项目路径：\path{/home/fanghaotian/Desktop/src/video-note-pipeline}
\item 报告工作目录：\path{/home/fanghaotian/Desktop/src/agent-basic-skill/reports/2026-03-26-bilibili-api-python-integration}
\item 安装版本：\texttt{bilibili-api-python==17.4.1}
\item 已验证请求客户端：\texttt{httpx}、\texttt{curl\_cffi}
\end{itemize}

\subsection{关键实验一览}

\begin{table}[H]
\centering
\begin{tabularx}{\textwidth}{p{3.4cm}p{2.2cm}Y}
\toprule
实验 & 结果 & 工程含义 \\
\midrule
cookie file 转 Credential & 成功 & 现有 \texttt{cookies-export} 结果可以直接复用，不需要新登录流程。 \\
Credential 有效性检查 & \texttt{check\_valid = true} & 当前浏览器导出的正式 cookie 文件可以作为 fallback 输入。 \\
Credential 刷新检查 & \texttt{check\_refresh = true} 但 \texttt{refresh()} 失败 & 缺 \texttt{ac\_time\_value}，不要依赖库内 refresh。 \\
匿名字幕探测 & 失败 & 仅靠匿名 API 不足以获得 Bilibili 官方字幕。 \\
带 cookies 的字幕探测 & 成功得到 \texttt{ai-zh} & 能补足当前 \texttt{yt-dlp} 可能漏掉的官方字幕路径。 \\
原始字幕 JSON 下载 & 成功，\texttt{body\_len = 387} & fallback 不止能知道“有字幕”，还足以生成标准 transcript。 \\
下载地址探测 & 公开和认证都可用，认证流更丰富 & 未来可作为视频画质探测的补充信息来源，但不建议本轮替换下载栈。 \\
\texttt{curl\_cffi} 实测 & 成功 & 可作为遇到风控时的增强模式。 \\
QR 登录验证 & 能生成二维码，未完成整链路登录 & 功能存在，但不建议暴露到当前产品交互。 \\
\bottomrule
\end{tabularx}
\caption{本地实验摘要}
\end{table}

\subsection{目标视频的最重要证据}

对用户关心的目标视频 \texttt{BV1DzALzQEHU}，本地已经拿到了以下关键结果：
\begin{itemize}
\item \texttt{public-vs-auth-behavior.json} 显示：公开可读 \texttt{title}、\texttt{pages\_count} 和 \texttt{cid}，但公开 \texttt{player\_info} 会因为缺 \texttt{SESSDATA} 失败。
\item 同一份 artifact 也显示：带认证后 \texttt{auth\_subtitle\_count = 1}，且第一条 \texttt{lan = ai-zh}。
\item \texttt{subtitle-metadata-and-json-summary.json} 显示：字幕索引对象中存在 \texttt{subtitle\_url} 与 \texttt{subtitle\_url\_v2}。
\item 原始字幕 JSON 样本大小约 69 KB，包含 387 段 \texttt{body[]}，足够归一化为项目 transcript。
\end{itemize}

\begin{evidencebox}{一个非常具体的判断}
这次实验说明，\texttt{bilibili-api-python} 的价值不在“猜语言标签”，而在于它能直接拿到 Bilibili 站内字幕索引和原始字幕 JSON 链接。也就是说，它最适合兜底“\texttt{yt-dlp} cookies 路径仍 miss，但 B 站官方字幕实际存在”的场景。
\end{evidencebox}

\section{应该如何接到当前项目里}

\subsection{推荐架构}

\begin{figure}[H]
\centering
\begin{tikzpicture}[
    node distance=10mm and 10mm,
    box/.style={draw, rounded corners, align=center, minimum height=9mm, text width=3.2cm, fill=blue!4},
    decision/.style={draw, diamond, aspect=2.1, align=center, inner sep=1pt, text width=2.9cm, fill=yellow!10},
    arrow/.style={-{Latex[length=2.2mm]}, thick}
]
\node[box] (probe) {现有 Bilibili\\subtitle probe};
\node[decision, below=of probe] (ytdlpok) {\texttt{yt-dlp}\\选中官方字幕?};
\node[box, left=of ytdlpok, xshift=-10mm] (done1) {沿用当前\\selected\_subtitle.json};
\node[decision, below=of ytdlpok] (cookieok) {本地 Bilibili\\cookie file 有效?};
\node[box, below=of cookieok] (bapi) {\texttt{bilibili-api-python}\\fallback};
\node[box, right=of ytdlpok, xshift=10mm] (asr) {继续走\\ASR fallback};
\node[box, below=of bapi] (normalize) {获取 \texttt{subtitle\_url}\\下载原始 JSON\\归一化 transcript};
\node[box, below=of normalize] (done2) {写入 probe artifact\\与 \texttt{selected\_subtitle.json}};

\draw[arrow] (probe) -- (ytdlpok);
\draw[arrow] (ytdlpok) -- node[above]{是} (done1);
\draw[arrow] (ytdlpok) -- node[right]{否} (cookieok);
\draw[arrow] (cookieok) -- node[right]{是} (bapi);
\draw[arrow] (cookieok) -- node[above]{否} (asr);
\draw[arrow] (bapi) -- (normalize);
\draw[arrow] (normalize) -- (done2);
\end{tikzpicture}
\caption{推荐的 Bilibili 字幕 fallback 架构}
\end{figure}

\subsection{依赖接入建议}

如果后续要写代码，建议不要直接把依赖塞进基础依赖组，而是新增一个 Bilibili 专用可选组。例如：

\begin{lstlisting}[caption={建议的 pyproject optional dependency 方向}]
[project.optional-dependencies]
bilibili-direct = [
  "bilibili-api-python>=17.4.1",
  "httpx>=0.28.0",
]
bilibili-direct-hardened = [
  "bilibili-api-python>=17.4.1",
  "curl_cffi>=0.14.0",
]
\end{lstlisting}

这样做的优点：
\begin{itemize}
\item 不影响 YouTube 用户或不需要 Bilibili fallback 的用户。
\item 把 GPLv3+ 依赖的进入点缩小到明确的 optional feature。
\item 让 \texttt{httpx} 与 \texttt{curl\_cffi} 的选择更清晰。
\end{itemize}

\subsection{同步/异步边界建议}

当前项目是同步 CLI，\texttt{bilibili-api-python} 是全异步库。建议新增一个很薄的内部模块，例如 \texttt{video\_note\_pipeline/adapters/bilibili\_api\_fallback.py}：
\begin{itemize}
\item 内部写一个 \texttt{async def \_probe\_subtitle\_index(...)} 异步函数。
\item 外层暴露同步函数，由 \texttt{asyncio.run(...)} 驱动。
\item 不要把异步对象泄漏到现有 command / adapter 边界之外。
\end{itemize}

官方文档也提供 \texttt{sync()} 包装器，但从项目可读性和边界控制上，显式的内部 async wrapper 更容易审计。

\subsection{推荐的落盘策略}

\begin{table}[H]
\centering
\begin{tabularx}{\textwidth}{p{4.3cm}p{3.5cm}Y}
\toprule
文件名建议 & 内容 & 为什么要存 \\
\midrule
\texttt{bilibili\_api\_probe.json} & 本次 fallback 的选择客户端、cid、错误、是否命中字幕 & 让 probe 失败原因可调试，不依赖运行日志。 \\
\texttt{bilibili\_player\_info.json} & \texttt{get\_player\_info()} 原始返回 & 记录 \texttt{need\_login\_subtitle}、\texttt{subtitle}、地区等站内信息。 \\
\texttt{bilibili\_subtitle\_index.json} & \texttt{get\_subtitle()} 返回的字幕索引对象 & 记录每条字幕轨的 \texttt{lan}、\texttt{lan\_doc}、\texttt{subtitle\_url}。 \\
\texttt{bilibili\_subtitle\_raw.<lan>.json} & 通过 \texttt{subtitle\_url} 下载的原始字幕 JSON & 保留源数据，便于后续归一化逻辑调整时重放。 \\
\texttt{selected\_subtitle.json} & 归一化后的最终字幕 payload & 与现有 pipeline 约定保持一致，让后续写作 / ASR fallback 无感知。 \\
\bottomrule
\end{tabularx}
\caption{推荐的 artifact 保存策略}
\end{table}

\subsection{字段归一化建议}

从当前样本看，最低可行归一化规则如下：
\begin{itemize}
\item \texttt{source}：建议写成 \texttt{official-cookies:bilibili-api-python}
\item \texttt{language}：保留原始 \texttt{lan}，例如 \texttt{ai-zh}
\item \texttt{normalized\_language}：映射到现有项目惯例，例如 \texttt{zh}
\item \texttt{segments[]}：由 \texttt{body[]} 映射，每段至少保留 \texttt{start}、\texttt{end}、\texttt{text}
\item \texttt{provenance}：补充 \texttt{subtitle\_url}、\texttt{cid}、\texttt{selected\_client}
\end{itemize}

\begin{evidencebox}{为什么原始 JSON 也要存}
因为 Bilibili 返回的不只是文本段，还有 \texttt{type}、\texttt{lang}、\texttt{music}、\texttt{location} 等辅助字段。今天我们可能只取 \texttt{from/to/content}，但后续若想做歌词过滤、音乐置信度过滤、字幕位置分析，必须保留原始数据。
\end{evidencebox}

\subsection{多 P 与 URL 语义}

官方文档说明 \texttt{page\_index} 与 \texttt{cid} 至少提供一个，且 \texttt{cid} 优先。对当前项目，建议：
\begin{itemize}
\item 先沿用现有 metadata 获取逻辑决定目标页。
\item 若 URL 已带 \texttt{p=}，则优先把它解析成对应的 \texttt{cid}。
\item fallback 只消费“已经确定好的页”，不要在 Bilibili 专用层重新发明页面选择逻辑。
\end{itemize}

\section{哪些地方最容易踩坑}

\begin{riskbox}{坑 1：把它当作“直接拿 transcript 的库”}
不是。它直接给我们的，是 Bilibili 原生字幕索引和原始 JSON 链接，而不是已经符合项目 schema 的 transcript。实现时必须再做一次 HTTP GET 和归一化。
\end{riskbox}

\begin{riskbox}{坑 2：以为现有 cookie file 可以自动 refresh}
不是。当前 cookies 文件能登录、能看字幕，但没有 \texttt{ac\_time\_value}。如果实现里误调用 \texttt{credential.refresh()}，会直接失败。
\end{riskbox}

\begin{riskbox}{坑 3：把登录能力也一起引进来}
技术上可以，产品上不应该。QR 登录、短信登录、密码登录会让当前项目多出一整套交互、状态轮询和错误恢复路径，明显超出这一轮目标。
\end{riskbox}

\begin{riskbox}{坑 4：忽略许可证}
官方仓库明确是 GPLv3+。这类问题不能等代码写完再看。
\end{riskbox}

\begin{riskbox}{坑 5：在主链路里过早替换 yt-dlp}
\texttt{yt-dlp} 仍然承担媒体下载、跨平台 cookie 导出与部分字幕能力；完全替换会把风险面扩大。当前最合理的方式是把 \texttt{bilibili-api-python} 放进 Bilibili 专用 fallback 层。
\end{riskbox}

\section{方案比较与推荐决策}

\begin{table}[H]
\centering
\begin{tabularx}{\textwidth}{p{3.2cm}p{2.5cm}p{2.3cm}p{2.0cm}Y}
\toprule
方案 & 官方字幕命中能力 & 依赖/复杂度 & 对用户交互影响 & 结论 \\
\midrule
只保留 \texttt{yt-dlp} & 中等，依赖 sidecar 结果 & 最低 & 无新增交互 & 已是基础能力，但无法覆盖所有 Bilibili 站内字幕场景。 \\
增加 \texttt{bilibili-api-python} fallback & 高，能直接读站内字幕索引与原始 JSON & 中等 & 无新增交互，可复用 cookie export & \textbf{推荐}。收益最大，工程代价可控。 \\
全面改为 \texttt{bilibili-api-python} 主链路 & 高，但需要补下载、事件循环、更多观测与许可证评估 & 最高 & 可能引入新调试与登录问题 & 现在不推荐。 \\
引入 \texttt{bilibili-api-python} 并开放 QR 登录 & 技术可行 & 很高 & 新增交互明显 & 不符合当前项目产品方向。 \\
\bottomrule
\end{tabularx}
\caption{Bilibili 字幕方案比较}
\end{table}

\begin{findingbox}{最终建议}
如果下一步是写 change / propose，我建议把目标描述成：
\begin{itemize}
\item ``为 \texttt{video-note-pipeline} 增加 Bilibili 专用 \texttt{bilibili-api-python} 字幕 fallback，复用现有 cookie export，不新增用户登录交互，并保留原始字幕 JSON artifact。''
\end{itemize}
而不要写成：
\begin{itemize}
\item ``把 Bilibili 全面迁移到 \texttt{bilibili-api-python}''，
\item ``增加二维码登录''，
\item ``做自动 cookies refresh''。
\end{itemize}
\end{findingbox}

\section{结论与后续问题}

\subsection{已经有充分支持的结论}

\begin{itemize}
\item \texttt{bilibili-api-python} 可以直接复用当前项目导出的 Bilibili cookies 文件。
\item 对目标视频 \texttt{BV1DzALzQEHU}，它在本机上已经成功取回 \texttt{ai-zh} 字幕索引与原始字幕 JSON。
\item 它适合放在 Bilibili fallback 层，而不是替换现有 \texttt{yt-dlp} 主链路。
\item 接入后不需要新增用户登录交互，继续使用 \texttt{video-note cookies-export bilibili} 即可。
\item 真正可用的持久化方案必须同时保存原始站内 JSON 和项目归一化 transcript。
\end{itemize}

\subsection{很可能成立，但仍建议在实现时再验证的判断}

\begin{itemize}
\item \texttt{httpx} 作为第一阶段默认客户端已经足够；只有在遇到 412 或更强风控时才需要切到 \texttt{curl\_cffi}。
\item \texttt{get\_download\_url()} 未来可能帮助视频画质探测，但这一轮没有必要动当前下载栈。
\item 当前 Bilibili fallback 的最佳触发条件，是 \texttt{yt-dlp} cookies 路径 miss 后再进入，而不是并行双探测。
\end{itemize}

\subsection{仍未知或值得单独验证的问题}

\begin{itemize}
\item 许可证约束在该项目的分发模式下是否可接受。
\item 多 P 视频、合集、番剧、课程等变体在 fallback 层是否需要单独处理。
\item 在更严格风控网络环境下，\texttt{httpx} 是否仍然足够稳定，还是需要默认 \texttt{curl\_cffi}。
\item \texttt{subtitle\_url\_v2} 是否比 \texttt{subtitle\_url} 更稳定，是否值得优先使用。
\end{itemize}

\appendix

\section{复现实验命令}

\begin{lstlisting}[caption={隔离环境中安装并验证 bilibili-api-python 的参考命令}]
cd /home/fanghaotian/Desktop/src/agent-basic-skill/reports/2026-03-26-bilibili-api-python-integration
uv venv artifacts/bapi-venv
uv pip install --python artifacts/bapi-venv/bin/python bilibili-api-python httpx
uv pip install --python artifacts/bapi-venv/bin/python curl_cffi
\end{lstlisting}

\begin{lstlisting}[caption={验证当前 cookie file 能否拿到目标视频字幕的参考命令}]
source artifacts/bapi-venv/bin/activate
python - <<'PY'
import asyncio
from pathlib import Path
from bilibili_api import Credential, video, select_client

def load_netscape_cookies(cookie_file: Path) -> dict[str, str]:
    cookies = {}
    for raw in cookie_file.read_text(encoding="utf-8").splitlines():
        if not raw or raw.startswith("#"):
            continue
        parts = raw.split("\t")
        if len(parts) >= 7:
            cookies[parts[5]] = parts[6]
    return cookies

async def main() -> None:
    select_client("httpx")
    credential = Credential.from_cookies(
        load_netscape_cookies(
            Path("/home/fanghaotian/.config/video-note-pipeline/cookies/bilibili.txt")
        )
    )
    v = video.Video(bvid="BV1DzALzQEHU", credential=credential)
    cid = await v.get_cid(0)
    subtitle = await v.get_subtitle(cid=cid)
    print(subtitle["subtitles"][0]["lan"])

asyncio.run(main())
PY
\end{lstlisting}

\section{来源清单}

\small
\begin{longtable}{p{1.2cm}p{1.6cm}p{4.0cm}p{6.4cm}}
\toprule
ID & 类型 & 来源 & 链接 / 路径 \\
\midrule
\endfirsthead
\toprule
ID & 类型 & 来源 & 链接 / 路径 \\
\midrule
\endhead
S1 & Primary & bilibili-api 官方 README & \url{https://github.com/Nemo2011/bilibili-api} \\
S2 & Primary & bilibili-api 官方文档入口 & \url{https://nemo2011.github.io/bilibili-api/} \\
S3 & Primary & PyPI 项目页 & \url{https://pypi.org/project/bilibili-api-python/} \\
S4 & Primary & 获取 Credential 文档 & \url{https://github.com/Nemo2011/bilibili-api/blob/main/docs/get-credential.md} \\
S5 & Primary & refresh cookies 文档 & \url{https://github.com/Nemo2011/bilibili-api/blob/main/docs/refresh_cookies.md} \\
S6 & Primary & 请求客户端文档 & \url{https://github.com/Nemo2011/bilibili-api/blob/main/docs/request_client.md} \\
S7 & Primary & 配置文档 & \url{https://github.com/Nemo2011/bilibili-api/blob/main/docs/configuration.md} \\
S8 & Primary & Video 模块文档 & \url{https://github.com/Nemo2011/bilibili-api/blob/main/docs/modules/video.md} \\
S9 & Primary & Login v2 模块文档 & \url{https://github.com/Nemo2011/bilibili-api/blob/main/docs/modules/login_v2.md} \\
S10 & Repo & 当前项目 Bilibili adapter & \path{/home/fanghaotian/Desktop/src/video-note-pipeline/src/video_note_pipeline/adapters/bilibili.py} \\
S11 & Repo & 当前 subtitle probe 逻辑 & \path{/home/fanghaotian/Desktop/src/video-note-pipeline/src/video_note_pipeline/probes/subtitles.py} \\
S12 & Repo & 当前依赖声明 & \path{/home/fanghaotian/Desktop/src/video-note-pipeline/pyproject.toml} \\
S13 & Experiment & Credential 可用性 & \path{/home/fanghaotian/Desktop/src/agent-basic-skill/reports/2026-03-26-bilibili-api-python-integration/artifacts/credential-check.json} \\
S14 & Experiment & refresh 失败摘要 & \path{/home/fanghaotian/Desktop/src/agent-basic-skill/reports/2026-03-26-bilibili-api-python-integration/artifacts/credential-refresh-attempt.json} \\
S15 & Experiment & 公开/认证差异 & \path{/home/fanghaotian/Desktop/src/agent-basic-skill/reports/2026-03-26-bilibili-api-python-integration/artifacts/public-vs-auth-behavior.json} \\
S16 & Experiment & 字幕元信息与 JSON 摘要 & \path{/home/fanghaotian/Desktop/src/agent-basic-skill/reports/2026-03-26-bilibili-api-python-integration/artifacts/subtitle-metadata-and-json-summary.json} \\
S17 & Experiment & 原始字幕 JSON & \path{/home/fanghaotian/Desktop/src/agent-basic-skill/reports/2026-03-26-bilibili-api-python-integration/artifacts/subtitle-json-sample.json} \\
S18 & Experiment & 下载地址能力摘要 & \path{/home/fanghaotian/Desktop/src/agent-basic-skill/reports/2026-03-26-bilibili-api-python-integration/artifacts/download-url-summary.json} \\
S19 & Experiment & QR 登录能力摘要 & \path{/home/fanghaotian/Desktop/src/agent-basic-skill/reports/2026-03-26-bilibili-api-python-integration/artifacts/qrcode-login-summary.json} \\
S20 & Experiment & \texttt{curl\_cffi} 字幕检查 & \path{/home/fanghaotian/Desktop/src/agent-basic-skill/reports/2026-03-26-bilibili-api-python-integration/artifacts/curl-cffi-subtitle-check.json} \\
\bottomrule
\end{longtable}

\normalsize

\end{document}
